\section{Additional Simulations}
\label{supp:simulations}

The simulation registry includes independent, periodic, noisy, mixture, missingness, and degradation families. Independent simulations check whether estimators approach the non-clustered target. Periodic simulations check whether the workflow detects clustered returns. Noisy simulations test whether finite perturbations mask recurrence, following the warning that observed noise can move finite samples toward apparent independence \citep{Faranda2013Noise}. Regime mixtures and preprocessing perturbations test whether clustering can be generated by analysis choices rather than physical degradation.

The generated simulation asset used in the main paper is the reproducibility-scale artifact. It is sufficient for build validation, but it is not a replacement for the full registered experiment matrix. The supplement therefore treats simulation results as evidence about the implemented pipeline and estimator sensitivity, not as an exhaustive finite-sample theory.

\subsection{Reference availability}
\label{supp:reference-availability}

Table~\ref{tab:estimator-reference-regimes} records, for every simulated family, how many configurations carry a reference extremal index, whether that reference is theoretical or numerical, and which clustering regime it falls in. Three families --- noisy periodic dynamics, persistent shift, and regime mixture --- have no established reference at any configuration: the two long-run estimates used to fix a numerical value disagree beyond the accepted tolerance, with discrepancies from \(0.15\) to \(0.60\) and no plateau in the block-based family. These are mixture and non-stationary processes for which a single extremal index is arguably not well defined, and they are reported as non-identifiable rather than assigned a value that one declustering convention happens to produce.

\input{../../reports/paper/latex/estimator_full_grid.tex}

The complete estimator grid, carrying all thirteen reported quantities for every configuration and estimator, is distributed as \texttt{estimator\_full\_grid.csv} in the reproducibility package. Every value in the main-paper table is recomputable from it.

\subsection{Sensitivity to perturbations}
\label{supp:simulation-sensitivity}

Figure~\ref{fig:supp-estimator-sensitivity} reports mean estimates under noise, missingness, and regime mixture. Added noise moves every estimator towards independence, which is the expected direction: perturbation breaks the deterministic recurrence that produces clustering, and the numerically established references for the noisy families sit between \(0.94\) and \(0.98\) accordingly. This is also why the high-replication grid, which runs at a single non-zero noise level, is reported in the main paper as evidence about near-independent processes.

\begin{figure}[h]
  \centering
  \includegraphics[width=\linewidth]{../../reports/paper/figures/estimator_sensitivity.png}
  \caption{Estimator sensitivity to noise scale, missing rate, and regime mixture. The regime-mixture panel contrasts families rather than sweeping a parameter, since mixture is a property of the process and not a tuning knob.}
  \label{fig:supp-estimator-sensitivity}
\end{figure}

\begin{figure}[h]
  \centering
  \includegraphics[width=0.78\linewidth]{../../reports/paper/figures/simulation_bias_rmse.png}
  \caption{Supplementary copy of the generated simulation RMSE artifact.}
  \label{fig:supp-simulation-rmse}
\end{figure}
